\documentclass[10pt,a4paper]{amsart}
\usepackage{amsmath,amssymb,amsthm,graphicx}
\usepackage{verbatim}
\newenvironment{reference}{\comment}{\endcomment}
%\newenvironment{reference}{}{}
\newenvironment{slogan}{\comment}{\endcomment}
\newenvironment{history}{\comment}{\endcomment}

% For commutative diagrams we use Xy-pic
\usepackage[all]{xy}

% We use 2cell for 2-commutative diagrams.
\xyoption{2cell}
% Theorem environments.
%
\theoremstyle{plain}
\newtheorem{theorem}[subsection]{Theorem}
\newtheorem{proposition}[subsection]{Proposition}
\newtheorem{lemma}[subsection]{Lemma}

\theoremstyle{definition}
\newtheorem{definition}[subsection]{Definition}
\newtheorem{example}[subsection]{Example}
\newtheorem{exercise}[subsection]{Exercise}
\newtheorem{situation}[subsection]{Situation}

\theoremstyle{remark}
\newtheorem{remark}[subsection]{Remark}
\newtheorem{remarks}[subsection]{Remarks}

\numberwithin{equation}{subsection}

% Macros
%
\def\lim{\mathop{\mathrm{lim}}\nolimits}
\def\colim{\mathop{\mathrm{colim}}\nolimits}
\def\Spec{\mathop{\mathrm{Spec}}}
\def\Hom{\mathop{\mathrm{Hom}}\nolimits}
\def\Ext{\mathop{\mathrm{Ext}}\nolimits}
\def\SheafHom{\mathop{\mathcal{H}\!\mathit{om}}\nolimits}
\def\SheafExt{\mathop{\mathcal{E}\!\mathit{xt}}\nolimits}
\def\Sch{\mathit{Sch}}
\def\Mor{\mathop{\mathrm{Mor}}\nolimits}
\def\Ob{\mathop{\mathrm{Ob}}\nolimits}
\def\Sh{\mathop{\mathit{Sh}}\nolimits}
\def\NL{\mathop{N\!L}\nolimits}
\def\CH{\mathop{\mathrm{CH}}\nolimits}
\def\proetale{{pro\text{-}\acute{e}tale}}
\def\etale{{\acute{e}tale}}
\def\QCoh{\mathit{QCoh}}
\def\Ker{\mathop{\mathrm{Ker}}}
\def\Im{\mathop{\mathrm{Im}}}
\def\Coker{\mathop{\mathrm{Coker}}}
\def\Coim{\mathop{\mathrm{Coim}}}

% Boxtimes
%
\DeclareMathSymbol{\boxtimes}{\mathbin}{AMSa}{"02}

%
% Macros for moduli stacks/spaces
%
\def\QCohstack{\mathcal{QC}\!\mathit{oh}}
\def\Cohstack{\mathcal{C}\!\mathit{oh}}
\def\Spacesstack{\mathcal{S}\!\mathit{paces}}
\def\Quotfunctor{\mathrm{Quot}}
\def\Hilbfunctor{\mathrm{Hilb}}
\def\Curvesstack{\mathcal{C}\!\mathit{urves}}
\def\Polarizedstack{\mathcal{P}\!\mathit{olarized}}
\def\Complexesstack{\mathcal{C}\!\mathit{omplexes}}
% \Pic is the operator that assigns to X its picard group, usage \Pic(X)
% \Picardstack_{X/B} denotes the Picard stack of X over B
% \Picardfunctor_{X/B} denotes the Picard functor of X over B
\def\Pic{\mathop{\mathrm{Pic}}\nolimits}
\def\Picardstack{\mathcal{P}\!\mathit{ic}}
\def\Picardfunctor{\mathrm{Pic}}
\def\Deformationcategory{\mathcal{D}\!\mathit{ef}}
\usepackage{cleveref}
\crefname{originalremark}{original remark}{original remarks}
\Crefname{originalremark}{Original remark}{Original remarks}
\makeatletter
\newlabel{varieties-remark-smooth-projective-compactification}{{0H1F}{source}{}{sourcepointer-1}{}}
\newlabel{varieties-remark-smooth-projective-compactification@cref}{{[originalremark][0][]0H1F}{[1][0][]source}{}{}{}}
\makeatother
\crefname{originallemma}{original lemma}{original lemmas}
\Crefname{originallemma}{Original lemma}{Original lemmas}
\makeatletter
\newlabel{lemma-cohomology-smooth-projective-curve}{{03RQ}{source}{}{sourcepointer-2}{}}
\newlabel{lemma-cohomology-smooth-projective-curve@cref}{{[originallemma][0][]03RQ}{[1][0][]source}{}{}{}}
\makeatother
\crefname{originallemma}{original lemma}{original lemmas}
\Crefname{originallemma}{Original lemma}{Original lemmas}
\makeatletter
\newlabel{lemma-describe-h1-mun}{{040Q}{source}{}{sourcepointer-3}{}}
\newlabel{lemma-describe-h1-mun@cref}{{[originallemma][0][]040Q}{[1][0][]source}{}{}{}}
\makeatother
\title{Cohomology of an affine smooth curve with roots of unity}
\author{The Stacks Project Authors}
\date{stacks-project:etale-cohomology.tex:local-snapshot-20261011; GNU FDL 1.2 or later}
\begin{document}
\maketitle
\noindent\textit{Editorial note.} All delivered mathematical text is the author's own, copied from the pinned Stacks Project chapter file etale-cohomology.tex; no statement or proof step was rewritten, reordered or removed.\par\medskip
\noindent\textit{Editorial note.} The proof is the author's complete computation for all three parts; the references to earlier discussions are pointers to notation already introduced, not omissions of arguments, and the exact sequence together with the rank count is written out.\par\medskip
\noindent\textit{Editorial note.} Every labelled reference occurring inside the counted spans is carried in assembly.external\_references as a non-load-bearing navigation pointer (same-chapter ones at their definition site, cross-chapter ones at the invocation inside the counted span), so no load-bearing mathematics is replaced by a stub. No bibliography citation occurs inside the counted spans.\par\medskip
\noindent\textit{Editorial note.} The checked chapter list for same-name collisions is the set of stacks chapters staged in this repository: chow, cohomology, divisors, etale-cohomology, flat, more-algebra, more-morphisms, morphisms, spaces-more-morphisms, varieties together with this card's chapter etale-cohomology.\par\medskip
\noindent\textit{Editorial note.} The preamble is the chapter's own native preamble excerpt already used for this source: the Xy-pic setup (native preamble lines 14-18) and the native theorem declarations with the mathematical macros (native preamble lines 154-217).\par\medskip
\noindent\textit{Allowed standard prerequisites (editorial).}\begin{itemize}
\item Cohomology of a smooth projective curve with roots of unity. Same chapter, Lemma lemma-cohomology-smooth-projective-curve. Supplies the rank count for the compactification and the divisibility discussion.
\item Description of the first cohomology in terms of triples. Same chapter, Lemma lemma-describe-h1-mun and the discussion preceding it. Supplies the description of the first cohomology group and of the isomorphism classes of triples.
\end{itemize}\medskip
\ifdefined\hypertarget\hypertarget{sourcepointer-1}{}\fi
\noindent\textit{Source reference.} Original remark 0H1F: sources/stacks-etale-cohomology/pinned-original-source.tex, source lines 11250--11250. Navigation pointer invoked at proof line 11250; the counted span contains the pointer only and the cited result is not reproduced inside it, so the pointer is registered as a non-load-bearing navigation stub in the established form for Stacks material. No bibliography citation is involved.\par\medskip
\ifdefined\hypertarget\hypertarget{sourcepointer-2}{}\fi
\noindent\textit{Source reference.} Original lemma 03RQ: sources/stacks-etale-cohomology/pinned-original-source.tex, source lines 11135--11135. Navigation pointer invoked at proof line 11234; the counted span contains the pointer only and the cited result is not reproduced inside it, so the pointer is registered as a non-load-bearing navigation stub in the established form for Stacks material. No bibliography citation is involved.\par\medskip
\ifdefined\hypertarget\hypertarget{sourcepointer-3}{}\fi
\noindent\textit{Source reference.} Original lemma 040Q: sources/stacks-etale-cohomology/pinned-original-source.tex, source lines 3057--3057. Navigation pointer invoked at proof line 11275; the counted span contains the pointer only and the cited result is not reproduced inside it, so the pointer is registered as a non-load-bearing navigation stub in the established form for Stacks material. No bibliography citation is involved.\par\medskip
\begin{lemma}
\label{lemma-vanishing-cohomology-mu-smooth-curve}
Let $X$ be an affine smooth curve over an algebraically closed field $k$ and
$n \in k^*$. Let $X \subset \overline{X}$ be a smooth projective
compactification
(Varieties, Remark \ref{varieties-remark-smooth-projective-compactification}).
Let $g$ be the genus of $\overline{X}$ and let $r$ be the number of
points of $\overline{X} \setminus X$. Then
\begin{enumerate}
\item $H_\etale^0(X, \mu_n) = \mu_n(k)$;
\item $H_\etale^1(X, \mu_n) \cong (\mathbf{Z}/n\mathbf{Z})^{2g+r-1}$, and
\item $H_\etale^q(X, \mu_n) = 0$ for all $q \geq 2$.
\end{enumerate}
\end{lemma}

\begin{proof}
Write $X = \overline{X} - \{ x_1, \ldots, x_r\}$. Then
$\Pic(X) = \Pic(\overline{X})/ R$, where $R$ is the subgroup generated by
$\mathcal{O}_{\overline{X}}(x_i)$, $1 \leq i \leq r$.
Since $r \geq 1$, we see that $\Pic^0(\overline{X}) \to \Pic(X)$
is surjective, hence $\Pic(X)$ is divisible (see discussion in proof
of Lemma \ref{lemma-cohomology-smooth-projective-curve}).
Applying the Kummer sequence, we get (1) and (3). For (2) we use
\begin{align*}
H_\etale^1(X, \mu_n)
& =
\{(\mathcal L, \alpha)\}/\cong \\
& =
(\{(\bar{\mathcal L}, D, \bar \alpha)\}/\cong)/\tilde{R}
\end{align*}
For the first equality sign and notation see Lemma \ref{lemma-describe-h1-mun}.
The triples $(\bar{\mathcal L}, D, \bar \alpha)$ in the second line
consist of an invertible module $\bar{\mathcal L}$ on $\overline{X}$
of degree $0$, a divisor $D$ on $\overline{X}$ of degree $0$ supported on
$\left\{x_1, \ldots, x_r\right\}$ and an isomorphism
$\bar{\alpha}: \bar{\mathcal L}^{\otimes n} \cong \mathcal{O}_{\bar{X}}(D)$.
Isomorphisms of triples are defined as in the discussion preceding
Lemma \ref{lemma-describe-h1-mun}.
Further $\tilde{R}$ is the subgroup generated by triples of the form
$(\mathcal{O}_{\overline{X}}(D'), n D', 1^{\otimes n})$
where $D'$ is supported on $\left\{x_1, \ldots, x_r\right\}$ and has degree 0.
The discussion above shows we have the second equality sign via
the map sending the triple $(\bar{\mathcal L},\ D,\ \bar \alpha)$
to the pair $(\bar{\mathcal L}|_X, \bar \alpha|_X)$. Thus,
we get an exact sequence
$$
0 \longrightarrow
H_\etale^1(\overline{X}, \mu_n) \longrightarrow
H_\etale^1(X, \mu_n) \longrightarrow
\bigoplus_{i = 1}^r \mathbf{Z}/n\mathbf{Z}
\xrightarrow{\ \sum\ }
\mathbf{Z}/n\mathbf{Z} \longrightarrow 0
$$
where the middle map sends the class of a triple
$(\bar{ \mathcal L}, D, \bar \alpha)$ with $D = \sum_{i = 1}^r a_i (x_i)$
to the $r$-tuple $(a_i)_{i = 1}^r$. It now suffices to use
Lemma \ref{lemma-cohomology-smooth-projective-curve}
to count ranks.
\end{proof}

\end{document}
